\chapter{임베딩과 포지셔널 인코딩}
\label{ch:embedding}

이 장에서는 토큰 ID를 밀집 벡터(dense vector)로 변환하는 \keyterm{임베딩(embedding)}과 시퀀스의 순서 정보를 모델에 주입하는 \keyterm{포지셔널 인코딩(positional encoding)}을 구현한다.

\section{왜 임베딩이 필요한가}

3장에서 토크나이저가 ``Hello''를 정수 ID 247로 바꿔주었다. 하지만 이 정수를 그대로 신경망에 넣으면 문제가 생긴다. 첫째, 정수 247과 248이 ``비슷한 단어''라는 보장이 없다. ``cat''이 247, ``dog''이 248일 수도 있고 ``zebra''가 248일 수도 있다. 둘째, 정수를 곱하거나 더하는 연산이 의미를 갖지 않는다. ``cat + 1 = dog''이 성립하지 않는다.

해결책은 각 토큰 ID를 \keyterm{밀집 벡터}로 변환하는 것이다. 어휘 크기가 $V$이고 임베딩 차원이 $d$일 때, 임베딩은 $V \times d$ 크기의 학습 가능한 행렬 $E$다. 토큰 ID $i$가 들어오면 $E$의 $i$번째 행을 가져온다. 이 행벡터가 그 토큰의 ``의미적 표현''이 된다.

\begin{infobox}[왜 ``밀집''인가?]
원-핫 인코딩(one-hot)은 어휘 크기 $V$의 희소 벡터로 토큰을 표현한다. 예를 들어 1만 어휘에 대해 ID 247은 246개의 0과 1개의 1, 나머지 0으로 이루어진 벡터다. 이를 $E$와 곱하면 결국 $E$의 247번째 행을 가져오는 것과 같다. 임베딩은 이 ``원-핫 $\times$ 행렬''을 lookup 연산 하나로 수행하는 효율적 구현이다.

학습이 진행되면 비슷한 의미의 단어가 벡터 공간에서 가까워진다. ``cat''과 ``dog''은 가깝고, ``cat''과 ``car''는 멀다. 이것이 임베딩의 마법이다: 학습 전에는 무작위 벡터였지만, 학습 후에는 의미를 담은 좌표가 된다.
\end{infobox}

\section{임베딩의 수학적 정의}

임베딩 행렬 $E \in \mathbb{R}^{V \times d}$가 있을 때, 토큰 ID $i$의 임베딩은 $E[i] \in \mathbb{R}^d$다. 이것은 단순히 $E$의 $i$번째 행을 선택하는 연산이며, 미분 가능하다 (backward 시 선택된 행만 그래디언트를 받는다).

배치 입력 $\mathbf{x} \in \mathbb{Z}^{B \times L}$ (배치 $B$, 시퀀스 길이 $L$)에 대해 임베딩 출력은 $\mathbf{X} \in \mathbb{R}^{B \times L \times d}$가 된다.

\section{Token Embedding 구현}

PyTorch의 \code{nn.Embedding}은 정확히 이 lookup 연산을 수행한다. 우리는 이를 얇게 래핑하여 초기화 로직을 추가한다.

\begin{lstlisting}[language=Python]
import math
import torch.nn as nn

class TokenEmbedding(nn.Module):
    def __init__(self, vocab_size: int, d_model: int, pad_id: int = 0):
        super().__init__()
        self.vocab_size = vocab_size
        self.d_model = d_model
        self.pad_id = pad_id
        self.embedding = nn.Embedding(vocab_size, d_model, padding_idx=pad_id)
        self._init_weights()

    def _init_weights(self):
        # 정규분포 N(0, 1/sqrt(d_model)) 초기화 (GPT-2 스타일)
        nn.init.normal_(self.embedding.weight, mean=0.0,
                        std=1.0 / math.sqrt(self.d_model))
        # 패딩 토큰은 0으로
        with torch.no_grad():
            self.embedding.weight[self.pad_id].fill_(0.0)

    def forward(self, token_ids: torch.Tensor) -> torch.Tensor:
        # [batch, seq] -> [batch, seq, d_model]
        return self.embedding(token_ids)
\end{lstlisting}

\begin{notebox}[초기화의 중요성]
딥러닝에서 가중치 초기화는 학습 성패를 좌우한다. 너무 크면 활성값이 폭발하고, 너무 작으면 그래디언트가 사라진다. GPT-2는 표준편차 $1/\sqrt{d_{model}}$의 정규분포를 사용하며, 이는 초기 임베딩 벡터의 norm이 대략 1이 되도록 하는 합리적 선택이다.

왜 하필 $1/\sqrt{d}$인가? $d$차원 벡터의 각 원소가 $\mathcal{N}(0, \sigma^2)$일 때, 벡터 norm의 기댓값은 $\sqrt{d} \cdot \sigma$다. $\sigma = 1/\sqrt{d}$면 norm이 1 근처가 된다. 이렇게 하면 임베딩 벡터의 스케일이 레이어 출력과 비슷하여 학습이 안정된다.
\end{notebox}

\subsection{padding\_idx의 역할}

\code{nn.Embedding}의 \code{padding\_idx} 매개변수는 패딩 토큰의 임베딩을 항상 0으로 유지한다. 이렇게 하면:
\begin{itemize}
  \item 패딩 위치의 출력이 0이 되어 어텐션 계산에 영향을 주지 않는다.
  \item backward 시 패딩 행의 그래디언트가 0이 되어 가중치가 갱신되지 않는다.
\end{itemize}

\section{포지셔널 인코딩}

트랜스포머는 어텐션만으로 시퀀스를 처리하는데, 어텐션 자체는 ``순서''라는 개념이 없다. 입력을 섞어도 출력은 (재배열된 채로) 같다. 이를 \keyterm{순서 불변성(permutation invariance)}이라고 한다. 자연어에서는 순서가 의미를 결정하므로(``개가 사람을 물었다'' vs ``사람이 개를 물었다''), 순서 정보를 명시적으로 주입해야 한다.

\subsection{직관: 왜 순서가 중요한가?}

어텐션은 ``각 토큰이 다른 토큰에 얼마나 주의를 기울일지''를 학습한다. 하지만 ``첫 번째 토큰''과 ``마지막 토큰''의 구분이 없다. 문장 ``나는 밥을 먹었다''에서 ``먹었다''가 주어 ``나는''과 연결되어야 하는데, 어텐션은 ``5번째 토큰이 1번째 토큰을 본다''는 정보 없이는 이 연결을 학습하기 어렵다.

포지셔널 인코딩은 각 위치에 고유한 ``위치 벡터''를 더해주어 모델이 위치를 구분할 수 있게 한다.

\subsection{Sinusoidal 위치 인코딩 (Vaswani et al. 2017)}

원본 트랜스포머 논문이 제안한 방식. 사인과 코사인 함수로 위치를 인코딩한다.

\begin{formula}[Sinusoidal Positional Encoding]
\[
PE_{(pos, 2i)} = \sin\left(\frac{pos}{10000^{2i/d_{model}}}\right)
\]
\[
PE_{(pos, 2i+1)} = \cos\left(\frac{pos}{10000^{2i/d_{model}}}\right)
\]
여기서 $pos$는 위치, $i$는 차원 인덱스. 짝수 차원은 sin, 홀수 차원은 cos.
\end{formula}

이 방식의 장점은 학습 가능한 파라미터가 없다는 것이다. 또한 상대적 위치에 대한 선형 관계가 성립하여 모델이 쉽게 학습할 수 있다. $\sin(\alpha + \beta) = \sin\alpha\cos\beta + \cos\alpha\sin\beta$이므로 위치 $pos + k$의 인코딩은 위치 $pos$의 인코딩의 선형 변환으로 표현된다.

\begin{lstlisting}[language=Python]
class PositionalEmbedding(nn.Module):
    def __init__(self, d_model: int, max_len: int = 512, mode: str = "learned"):
        super().__init__()
        self.mode = mode
        if mode == "sinusoidal":
            pe = torch.zeros(max_len, d_model)
            position = torch.arange(0, max_len).unsqueeze(1).float()
            div_term = torch.exp(
                torch.arange(0, d_model, 2).float() * (-math.log(10000.0) / d_model)
            )
            pe[:, 0::2] = torch.sin(position * div_term)
            pe[:, 1::2] = torch.cos(position * div_term)
            self.register_buffer("pe", pe.unsqueeze(0))  # [1, max_len, d_model]
        else:  # learned
            self.pe = nn.Embedding(max_len, d_model)
            nn.init.normal_(self.pe.weight, mean=0.0, std=1.0/math.sqrt(d_model))

    def forward(self, seq_len: int) -> torch.Tensor:
        # [1, seq_len, d_model] 반환
        if seq_len > self.max_len:
            raise ValueError(f"Sequence length {seq_len} exceeds max_len {self.max_len}")
        if self.mode == "sinusoidal":
            return self.pe[:, :seq_len, :]
        else:
            positions = torch.arange(seq_len, device=self.pe.weight.device)
            return self.pe(positions).unsqueeze(0)
\end{lstlisting}

\begin{figure}[htbp]
\centering
\begin{tikzpicture}[
  scale=0.7,
  every node/.style={font=\small\sffamily}
]
  % sin 파형
  \draw[inkblue, line width=0.5pt, ->] (0, 0) -- (8, 0) node[right] {pos};
  \draw[inkblue, line width=0.5pt, ->] (0, -1.2) -- (0, 1.2) node[above] {value};
  \draw[inkred, line width=1pt, domain=0:7, samples=50] plot (\x, {sin(\x*50)});
  \node[inkred, font=\scriptsize] at (7, 0.8) {dim 0 (sin)};
  \draw[inkblue, line width=1pt, domain=0:7, samples=50] plot (\x, {cos(\x*50)});
  \node[inkblue, font=\scriptsize] at (7, -0.8) {dim 1 (cos)};
\end{tikzpicture}
\caption{Sinusoidal 위치 인코딩의 처음 두 차원. sin과 cos이 위치에 따라 주기적으로 변한다.}
\label{fig:sinusoidal}
\end{figure}

\subsection{학습 가능한 위치 임베딩 (GPT-2 스타일)}

GPT-2, GPT-3 등은 위치마다 학습 가능한 임베딩 벡터를 사용한다. 이 방식은 충분한 데이터가 있을 때 더 유연하지만, 학습 시 본 적 없는 길이(예: 학습 시 1024였는데 추론 시 2048)에는 대응하기 어렵다.

\begin{notebox}[Sinusoidal vs Learned]
\begin{itemize}
  \item \textbf{Sinusoidal}: 파라미터 없음, extrapolation 가능 (학습 시 본 길이보다 긴 시퀀스에도 대응). 하지만 실제 성능은 학습 가능 방식보다 약간 떨어지는 경향.
  \item \textbf{Learned}: 성능 약간 우수, 하지만 고정 길이. extrapolation 어려움.
\end{itemize}
최신 모델들은 RoPE(Rotary Position Embedding)를 사용하여 두 방식의 장점을 결합한다. 2권 5장에서 다룬다.
\end{notebox}

\section{임베딩 합성}

토큰 임베딩과 위치 임베딩을 더해서 최종 입력 표현을 만든다. Vaswani et al.은 임베딩에 $\sqrt{d_{model}}$을 곱하여 스케일을 맞추는 것을 제안했는데, 이는 위치 임베딩과 스케일을 비슷하게 유지하기 위함이다.

\begin{lstlisting}[language=Python]
class EmbeddingStack(nn.Module):
    def __init__(self, vocab_size, d_model, max_len=512,
                 pad_id=0, pos_mode="learned", dropout=0.1):
        super().__init__()
        self.token_emb = TokenEmbedding(vocab_size, d_model, pad_id=pad_id)
        self.pos_emb = PositionalEmbedding(d_model, max_len, mode=pos_mode)
        self.dropout = nn.Dropout(dropout)
        self.scale = math.sqrt(d_model)

    def forward(self, token_ids: torch.Tensor) -> torch.Tensor:
        batch, seq = token_ids.shape
        tok = self.token_emb(token_ids) * self.scale  # 스케일링
        pos = self.pos_emb(seq)
        return self.dropout(tok + pos)
\end{lstlisting}

\begin{figure}[htbp]
\centering
\begin{tikzpicture}[
  every node/.style={font=\small\sffamily},
  box/.style={draw=inkblue, fill=inkblue!8, rounded corners=2pt, minimum width=2.5cm, minimum height=0.7cm},
  sum/.style={draw=inkred, fill=inkred!10, circle, minimum size=0.6cm}
]
  \node[box] (ids) at (0, 0) {token IDs};
  \node[box, below=0.6cm of ids] (tok) {Token Embedding};
  \node[box, right=2cm of ids] (pos) {Position};
  \node[box, below=0.6cm of pos] (posemb) {Positional Embedding};
  \node[sum, below=0.8cm of $(tok)!0.5!(posemb)$] (add) {$+$};
  \node[box, below=0.6cm of add] (out) {Dropout $\to$ output};
  \draw[-{Stealth[length=4pt]}, inkblue] (ids) -- (tok);
  \draw[-{Stealth[length=4pt]}, inkblue] (pos) -- (posemb);
  \draw[-{Stealth[length=4pt]}, inkblue] (tok) -- (add);
  \draw[-{Stealth[length=4pt]}, inkblue] (posemb) -- (add);
  \draw[-{Stealth[length=4pt]}, inkblue] (add) -- (out);
\end{tikzpicture}
\caption{EmbeddingStack 구조. 토큰 임베딩과 위치 임베딩을 더한 후 드롭아웃을 적용한다.}
\label{fig:embedding-stack}
\end{figure}

\subsection{왜 스케일링이 필요한가?}

토큰 임베딩은 $\mathcal{N}(0, 1/d)$로 초기화되어 있어 원소의 분산이 $1/d$다. 위치 임베딩 (sinusoidal)의 원소는 $[-1, 1]$ 범위로 분산이 약 $1/2$이다. 스케일링 없이 더하면 위치 정보가 토큰 정보를 압도한다.

$\sqrt{d}$를 곱하면 토큰 임베딩의 분산이 $1$이 되어 위치 임베딩과 균형을 이룬다.

\section{Layer Normalization 소개}

트랜스포머에서 학습을 안정화하기 위해 \keyterm{레이어 정규화(Layer Normalization)}를 사용한다. 다음 장의 어텐션에서도 등장하므로 여기서 소개한다.

\begin{formula}[Layer Normalization]
\[
\text{LN}(x) = \gamma \cdot \frac{x - \mu}{\sqrt{\sigma^2 + \epsilon}} + \beta
\]
여기서 $\mu, \sigma^2$은 마지막 차원을 따라 계산한 평균과 분산. $\gamma, \beta$는 학습 가능한 파라미터.
\end{formula}

PyTorch에서는 \code{nn.LayerNorm(d\_model)} 한 줄로 사용한다. 2권에서는 LayerNorm의 변형인 RMSNorm도 다룬다.

\begin{infobox}[LayerNorm vs BatchNorm]
BatchNorm은 배치 차원을 따라 정규화하여 학습과 추론 시 통계가 다른 문제가 있다. LayerNorm은 각 샘플 내부에서 정규화하여 배치 크기에 무관하게 동작한다. 이것이 LLM에서 LayerNorm을 쓰는 이유다: 배치 크기 1인 추론에서도 안정적으로 동작한다.
\end{infobox}

\section{테스트}

임베딩 모듈의 핵심 테스트는 출력 shape 확인과 그래디언트 흐름 확인이다.

\begin{lstlisting}[language=Python]
def test_token_embedding_shape():
    emb = TokenEmbedding(vocab_size=100, d_model=32)
    ids = torch.randint(0, 100, (2, 8))  # [batch=2, seq=8]
    out = emb(ids)
    assert out.shape == (2, 8, 32)

def test_positional_sinusoidal():
    pe = PositionalEmbedding(d_model=32, max_len=64, mode="sinusoidal")
    out = pe(seq_len=16)
    assert out.shape == (1, 16, 32)

def test_embedding_stack():
    stack = EmbeddingStack(vocab_size=100, d_model=32, max_len=64, dropout=0.0)
    ids = torch.randint(0, 100, (2, 8))
    out = stack(ids)
    assert out.shape == (2, 8, 32)

def test_positional_exceeds_max_len_raises():
    pe = PositionalEmbedding(d_model=32, max_len=32, mode="learned")
    with pytest.raises(ValueError):
        pe(seq_len=64)  # max_len 초과
\end{lstlisting}

\section{실습: 임베딩 시각화}

\code{scripts/tiny\_train.py}에서 학습된 임베딩을 PCA로 2차원에 투영하면 단어들이 어떻게 분포하는지 볼 수 있다. 충분한 학습 후에는 비슷한 단어들이 가깝게 모이는 것을 확인할 수 있다.

\section{요약}

\begin{itemize}
  \item 임베딩은 토큰 ID를 밀집 벡터로 변환하는 학습 가능한 lookup 테이블이다.
  \item 초기화는 $N(0, 1/\sqrt{d_{model}})$을 사용하고, 패딩 토큰은 0으로 설정한다.
  \item 포지셔널 인코딩은 트랜스포머에 순서 정보를 주입한다. sinusoidal과 learned 두 방식이 있다.
  \item 토큰 임베딩과 위치 임베딩을 더하고 드롭아웃을 적용하여 최종 입력 표현을 만든다.
  \item $\sqrt{d_{model}}$ 스케일링으로 두 임베딩의 분산을 균형 맞춘다.
  \item LayerNorm은 배치 크기에 무관한 정규화로 LLM 학습 안정화에 필수.
\end{itemize}

\section{연습 문제}

\begin{enumerate}
  \item 임베딩 차원 $d$가 너무 작으면 (예: 8) 어떤 문제가 생기는가? 너무 크면 (예: 4096) 어떤 문제가 생기는가?
  \item Sinusoidal 위치 인코딩에서 \code{10000}이라는 상수를 100으로 바꾸면 어떤 일이 일어나는가?
  \item 토큰 임베딩에 $\sqrt{d_{model}}$을 곱하지 않으면 학습이 어떻게 달라지는가?
  \item 학습 가능한 위치 임베딩의 \code{max\_len}을 1024로 설정하고 학습한 후, 추론 시 2048 길이 시퀀스를 넣으면 어떻게 되는가?
  \item LayerNorm의 $\epsilon$ 값을 $10^{-5}$에서 $10^{-1}$로 바꾸면 출력이 어떻게 변하는가?
\end{enumerate}

다음 장에서는 이 임베딩을 입력으로 받아 ``각 토큰이 다른 토큰에 얼마나 주의를 기울일지''를 학습하는 어텐션 메커니즘을 구현한다.
